% !TeX program = xelatex
\documentclass[12pt,a4paper]{ctexart}
\usepackage{amsmath,amssymb,bm}
\usepackage{booktabs,array,longtable}
\usepackage{graphicx}
\usepackage[margin=2.5cm]{geometry}
\usepackage{caption}
\captionsetup{font=small,labelsep=quad}
\usepackage{listings}
\renewcommand{\lstlistingname}{代码}
\usepackage{float}
\newfontfamily\codefont{JuliaMono-Regular.ttf}[
  Path=C:/texlive/2026/texmf-dist/fonts/truetype/public/juliamono/, Scale=0.875,
  BoldFont=JuliaMono-Bold.ttf, ItalicFont=JuliaMono-RegularItalic.ttf,
  BoldItalicFont=JuliaMono-BoldItalic.ttf]
\lstset{basicstyle=\ttfamily\codefont\small,breaklines=true,columns=flexible,frame=single,numbers=left,numberstyle=\tiny,
        showstringspaces=false,xleftmargin=15pt,framexleftmargin=11.6pt,xrightmargin=3.4pt,numbersep=5pt}
\usepackage[colorlinks=true,linkcolor=black,citecolor=black,urlcolor=black]{hyperref}
\graphicspath{{../05-图表/}}
\numberwithin{equation}{section}
\renewcommand{\arraystretch}{1.25}

\pagestyle{plain}

\linespread{1.0}

\ctexset{
  section = {name = {,、}, number = \chinese{section},
             format = \centering\songti\bfseries\zihao{4}},
  subsection = {number = \arabic{section}.\arabic{subsection},
                format = \songti\bfseries\zihao{-4}},
  subsubsection = {number = \arabic{section}.\arabic{subsection}.\arabic{subsubsection},
                   format = \songti\bfseries\zihao{-4}},
}

\renewenvironment{abstract}{%
  \normalsize
  \begin{center}{\songti\bfseries\zihao{4}\abstractname}\end{center}
  \par
}{\par}

\newcommand{\figplaceholder}[1]{\fbox{\parbox[c][5.5cm][c]{0.86\linewidth}{\centering 图位：#1}}}
\newcommand{\includefig}[3][0.86\linewidth]{\IfFileExists{../05-图表/#2}{\includegraphics[width=#1]{#2}}{\figplaceholder{#3}}}

\makeatletter
\renewcommand\@maketitle{%
  \newpage
  \null
  \vskip 2em%
  \begin{center}%
  \let\footnote\thanks
    {\LARGE\@title\par}%
    \vskip 1em%
    {\large\@date}%
  \end{center}%
  \par
  \vskip 1.5em}
\makeatother

\title{\textbf{中药材烘干热湿耦合建模、守恒求解与模型检验}\\[2pt]
——基于质量坐标有限体积的四问统一求解}
\author{}
\date{}

\begin{document}
\maketitle

\begin{abstract}

中药材干燥是产地加工的关键环节：水分偏高易霉变，干燥不足或过度都会影响品质。本文以长 25 cm、半径 2 cm 的圆柱药柱为对象，初始干基含水率 2.55 kg/kg、初温 28
${}^\circ$C，研究预热段与全过程的温度场与水分浓度场演变，并确定固定几何与计入失水收缩两种情形下的达标时间，为烘干时长的控制提供依据。

针对问题 1（0--1800 s 预热段），含水率变化很小、热湿近似解耦，材料参数按附录 2 取常数，热方程退化为线性圆柱 Robin 问题；由此写出圆柱 Robin 特征展开与时变环境 Duhamel
卷积的解析谱解，再与守恒型有限体积数值解逐点对照，最大偏差$5.9\times10^{-7}\,^{\circ}$C，1800 s 末全场温差收敛到约 $1\,^{\circ}$C。解析解不含离散误差，两者一致说明所报温度场是模型的解，空
间离散、边界处理与时间步进没有引入可观误差。

针对问题 2，材料参数按附录 3 取用，保留温度与水分浓度的双向耦合，用质量坐标下的守恒型有限体积格式求解：表面取半单元串联阻力，时间推进用后向
欧拉加局部外推，非线性项用 Picard--Newton 交替迭代，步长自适应。判据$Bi_m\in[1.1805,47.7480]$ 表明内部扩散主导，把含水率作集总或平均处理会把达标时间低估 46.03 h（$-80\%$），与题目要求的空间分布不
符；因此保留全场分布，给出 72 h 内各时刻的温度与水分浓度场。

针对问题 3，在问题 2 的模型上加入达标判据（各处水分浓度首次低于0.15 kg/kg 干基），由数值扫描确认中心为最慢点，按穿越区间及其右侧安全端点报告达标时刻：$\bm{t_f=57.1799}$ h（Richardson
外推 57.1636 h）。该值远大于常系数参考时标15.70 h，量级与参数一致；其中对边界汇取值最敏感，$C_e=0.1366$ 时增大 22.6 h，已检验因素的累计偏移量级约 1.71 h。

针对问题 4，几何按仿射收缩 $r=R(t)\sqrt q$ 演化，方程与固定域同形；为确定$R(t)$，把失水收缩、毛细-膨压塌陷、结皮-成孔停滞三项机制写成半经验关系式，由湿度场逐步预测半
径，预测不使用附件 2 的数据，附件 2 实测半径仅作交叉验证（RMSE 0.047 mm、偏差 $-0.56\%$）。得 $\bm{t_f=50.5369}$ h，比问题 3 快 $11.6\%$：附录 4 的材料参数使固定半径下的烘干时间放大 $2.26$ 倍（$+78.57$ h），半径收缩（$R^2$ 降
64\%）则加速$60.9\%$，两者相抵后反而更快；耦合主解在满 72 h 真实轨迹上的全程 RMSE 0.088 mm。

\smallskip\noindent\textbf{关键词}：热湿耦合\enspace 质量坐标\enspace 守恒有限体积\enspace 无量纲判据\enspace 模型检验\enspace 中药材干燥
\end{abstract}

\clearpage

\section{问题重述}

\subsection{问题背景}
中药材的干燥是药材加工的关键环节：水分过高会导致霉变，干燥不足或过度都会影响品质。
本题药材为圆柱形（长 $L_0=25$ cm、半径 $R_0=2$ cm），初始温度 $28\,^{\circ}$C、
初始干基含水率 2.55 kg/kg，置于烘干房中干燥；烘房空气温度与空气水分浓度由附件 1
给出（0--14400 s、60 s 一步），药材失水时外半径随之收缩（附件 2 给出 0--72 h、
1800 s 一步的半径，由 2.0 cm 缩至 1.198 cm 后进入近平台段）。烘干达标指药材各处
水分浓度（干基含水率）均首次低于 0.15 kg/kg。

\subsection{问题提出}
四个问题依次为：

\noindent \textbf{问题 1}：预热平衡阶段（0--1800 s）药材内温度与水分浓度的时空分布，
材料参数取附录 2，按题目所给的表格格式列出；

\noindent \textbf{问题 2}：整个烘干过程的温度与水分浓度分布情况，材料参数取附录 3；

\noindent \textbf{问题 3}：固定几何下首次全部达标的烘干时间，结果按表 5 的格式给出；

\noindent \textbf{问题 4}：计入失水收缩后的达标时间，外半径取附件 2、材料参数取附录 4，
结果按表 6 的格式给出。

四个结果分别写入 result1.xlsx 至 result4.xlsx。

总体思路是建立统一模型，四问共用\textbf{同一套热湿耦合控制方程}，差别只在取法：
材料参数来源、热湿耦合方式、几何是否演化以及结果形式四项不同。其中问题 1 是
问题 2 的解耦特例，问题 3 在问题 2 的基础上增加阈值判据，问题 4 是移动域的推广；
对收缩几何，先按有物理含义的半经验关系式由内而外地预测半径，附件 2 只作验证数据，
再用力学驱动扩展模型讨论收缩的机制边界。几何模型与边界条件见图~\ref{fig:几何示意}。

\begin{figure}[htbp]
\centering
\includefig[\linewidth]{fig17_几何模型与边界条件示意图.pdf}{几何模型与边界条件示意}
\caption{几何模型与边界条件。长 25 cm、半径 2 cm 的圆柱药柱，
初始 $U_0=2.55$ kg/kg、$T_0=301.15$ K；柱面与环境 $T_a(t)$、$C_{\mathrm{env}}(t)$
经对流换热 $h=25$ W/(m$^2\cdot$K) 与表面水分交换 $h_m=8\times10^{-7}$ m/s 耦合；
轴心对称、端面忽略（一维径向假设，见第三章）}
\label{fig:几何示意}
\end{figure}

\section{问题分析}

\subsection{四个问题的统一结构与难点}
四个问题（以下简称四问）的物理过程相同，都包含内部导热与导湿、表面对流边界以及材料参数的强非线性。
\textbf{(i)} 三套附录材料参数必须按问题分别使用，附录 2/3/4 的 $\rho,c_p,k,D$ 互不相同，
不得把三套参数调和、平均或取为常数。例如附录 4 与附录 3 的扩散系数比
$0.175\,\mathrm e^{0.15/U}$ 在高含水端约为 $1/5.4$（图~\ref{fig:扩散系数}），
\textbf{(ii)} $D\propto\mathrm e^{-\alpha/C}\mathrm e^{-3850/T}$ 对局部含水率与温度
强非线性，尾段 $D$ 最多缩小到原先的 16.8 倍，形成低扩散率的“干壳”，决定达标时间的尾段行为；
\textbf{(iii)} 环境数据只有 4 h，半径数据有 72 h，而烘干需 50 h 以上，边界外推方式与
收缩关系式的取法都会改变结论的量级。

\begin{figure}[htbp]
\centering
\includefig[\linewidth]{fig11_扩散系数与时间尺度.pdf}{扩散系数与特征时间尺度}
\caption{扩散系数非线性与特征时间尺度：(a) 三套附录材料参数的
$D$–$C$ 半对数曲线——附录 2 只与 $C$ 有关、不含 $T$，附录 3/4 为 $D(C,T)$；附录 4 全程低于附录 3，
尾段 $C\to0.15$ 时 $D$ 缩小 16.8 倍，该倍率按附录 3 的材料参数公式算出；
(b) 各参考时标：常系数首模参考 15.7021 h、膜界 11.1720 h、问题 4 瞬时首模时间积分
19.016 h、固定 $R_0$ 版 37.61 h、$R_{\min}$ 常值版 15.95 h——达标时间的
两个主值 57.18 h 与 50.54 h 均大于上述各时标}
\label{fig:扩散系数}
\end{figure}

\subsection{各问分析与建模路线}
\textbf{问题 1（预热段场）}：0--1800 s 含水率变化很小，热湿近似解耦，$D$ 不含 $T$、
按附录 2 取常数材料参数，热方程退化为线性圆柱 Robin 问题，解析谱解可直接写出，
其形式为圆柱 Robin 特征展开与时变环境的 Duhamel 卷积。把解析解与数值解逐点对照，
两者吻合到 $5.9\times10^{-7}\,^{\circ}$C，说明数值解求的正是这组方程，问题 1 的温度场
不是求解过程带进来的偏差；两边用的是同一套材料参数，所以这一比较不能说明材料参数的
取值本身是否正确（图~\ref{fig:解析验证}）。

\textbf{问题 2/3（全程场与达标时间）}：必须保留双向耦合（$D$ 含 $C,T$）与空间分布。
判据 $Bi_m\in[1.1805,47.7480]$ 表明内部扩散主导，任何“集总/平均含水率”的处理
都会把达标时间低估 46.03 h（$-80\%$），无法得到本题的达标时间。求解路线是质量坐标
守恒有限体积加自适应时间积分，并用解析退化、守恒恒等式与网格收敛一并检验；
达标判据用未舍入值在重建场上判定，中心为最慢点，已由逐步数值扫描确认。

\textbf{问题 4（收缩下达标时间）}：几何演化使方程成为移动域问题。在材料坐标下取
仿射收缩 $r=R(t)\sqrt q$，方程保持同形且无伪对流项；半径 $R(t)$ 的给定方式直接决定
问题 4 的结果。本文把“失水收缩、毛细-膨压塌陷、结皮-成孔停滞”三项有物理含义的机制
写成半经验关系式，由湿度场逐步\textbf{内生预测} $R(t)$，半径不使用附件 2 的数据；
附件 2 实测半径降为验证数据。另用力学驱动扩展模型讨论收缩的机制结构与精度边界，
该模型由有限应变径向平衡与可标定保水本构构成，见第 5.5 节。

\subsection{输入数据与预处理}
附件 1 的环境时程共 241 点，用分段线性插值取值，4 h 后采用末值保持外推，
外推值取 $T_a=50.165\,^{\circ}$C、$C_{\mathrm{env}}=0.04986$；其他取值方式作为对照，
已纳入敏感性分析。附件 2 半径用 PCHIP 保形插值，插值结果半径为正且单调不增，
72 h 后末值延拓，数据与插值行为见图~\ref{fig:输入数据}。

\begin{figure}[htbp]
\centering
\includefig[\linewidth]{fig1_输入数据.pdf}{输入数据}
\caption{输入数据：附件 1 给出 0--4 h、每 60 s 一步的烘房温度与空气水分浓度，
4 h 后按末值保持外推；附件 2 给出 0--72 h 的外半径，呈急缩、缓缩、平台三段结构，
约 21 h 起进入近平台段}
\label{fig:输入数据}
\end{figure}

\section{模型假设}
\begin{enumerate}
  \item \textbf{一维径向、轴对称，忽略端面与轴向}。依据：端面与侧面的面积比 $R/L=8\%$、
        轴向扩散路径 $L/2=12.5$ cm 是径向的 6.25 倍，扩散时间尺度为径向的
        $6.25^2\approx39$ 倍。
  \item \textbf{湿边界取本文设定的等效边界条件}：把附件 1 的 $C_{\mathrm{env}}(t)$
        直接当作与 $h_m$ 配套的 Robin 驱动势 $-D\partial_rU=h_m(U_s-C_{\mathrm{env}})$。
        依据：空气与固体含水率基准不同，严格的边界条件需借助吸附等温线或水活度关系；
        本文结论对这一取值最敏感，其影响由 $C_e=0.1366$ 的敏感性 $+22.6261$ h 给出，
        见第七章。
  \item \textbf{热方程取无潜热的有效显热近似}。依据：潜热通量与显热通量同量级，
        通量比为 2.06--2.49，用第一问 1800 s 点的表面量复算，通量比为 5.54；
        作为对照的潜热情景下 $t_f$ 由 57.1799 h 变为 62.5694 h（$+5.3889$ h），
        见 5.3 节。
  \item \textbf{材料参数逐点、按局部瞬时值取值}：$\rho,c_p,k,D$ 按附录公式在局部 $C$、
        局部 $T$（开尔文）取值；$h=25$ W/(m$^2\cdot$K)、$h_m=8\times10^{-7}$ m/s
        自附录 2 延用到问题 2--4。
  \item \textbf{环境外推}：$t>14400$ s 取末值保持。对照取法有线性外推与末 1 h 均值，
        差异在 $+1.62$ h 以内，已纳入敏感性分析。
  \item \textbf{几何}：长度不变，内部取等比仿射收缩 $r=R(t)\sqrt q$；$R(t)$ 由半经验
        关系式逐步预测，72 h 后保持 1.198 cm。
  \item \textbf{$\rho(C)$ 不用于质量—体积换算}：$\rho$ 只用于有效体积热容
        $a_{\mathrm{eff}}=\rho c_p$ 与通量算子 $\rho_{sv}=\rho/(1+C)$。依据：附录 3/4 的
        $\rho(C)$ 与附件 2 的 $R(t)$ 在干固体质量守恒下不相容，要求 $R\ge1.301$ cm
        而数据终值为 1.198 cm。
\end{enumerate}

\section{符号说明}
\begin{table}[H]
\centering
\caption{主要符号}
\label{tab:符号}
\small
\begin{tabular}{lll}
\toprule
符号 & 含义 & 单位 \\
\midrule
$q$ & 材料坐标（干物质标签），$q=(X/R_0)^2$ & — \\
$U(q,t)$ & 干基含水率（水分浓度） & kg/kg \\
$T(q,t)$ & 温度（开尔文） & K \\
$D(C,T)$ & 扩散系数（附录 2/3/4 公式） & m$^2$/s \\
$\rho,\ c_p,\ k$ & 密度、比热容、导热系数（逐点取值） & kg/m$^3$、J/(kg$\cdot$K)、W/(m$\cdot$K) \\
$a_{\mathrm{eff}}$ & 有效体积热容 $\rho c_p$ & J/(m$^3\cdot$K) \\
$h,\ h_m$ & 表面对流换热系数、表面传质系数 & W/(m$^2\cdot$K)、m/s \\
$T_a(t),\ C_{\mathrm{env}}(t)$ & 环境温度、环境水分浓度（附件 1） & K、kg/kg \\
$R(t)$ & 药柱外半径（问题 4 中为待求量） & m \\
$t_f$ & 达标时间（$\max_q U<0.15$ 首次成立） & h \\
$Le,\ Bi_m,\ Bi_h$ & 无量纲判据：$\alpha/D$、$h_mR/D$、$hR/k$ & — \\
$J$ & 体积变形梯度（力学扩展） & — \\
$S$ & 液相饱和度（力学扩展） & — \\
$p_c,\ \bar p_c$ & 毛细压及其储能平均（力学扩展） & （模型单位） \\
\bottomrule
\end{tabular}
\end{table}

\section{模型建立与求解}

\subsection{统一控制方程与质量坐标}
\textbf{守恒律与材料坐标}：从干固体守恒与水分守恒出发，
以干物质标签 $q\in[0,1]$ 作为材料坐标，其中 $q=(X/R_0)^2$、$X$ 为初始物质点，
描述径向一维热湿耦合过程；毛细多孔介质热湿传递的经典框架见~\cite{luikov}。
当干物质密度沿空间均匀、且内部取仿射收缩 $r=R(t)\sqrt q$ 时，
方程与固定域同形、无伪对流项，推导见附录 A：
\begin{equation}
U_t=\partial_q\!\left(\frac{4qD}{R(t)^2}U_q\right),\qquad
a_{\mathrm{eff}}\,T_t=\partial_q\!\left(\frac{4qk}{R(t)^2}T_q\right).
\label{eq:质量坐标}
\end{equation}
边界为移动表面 $q=1$ 上的 Robin 条件，形式不变，
其中跳条件里的几何运动项恒等相消：
\begin{equation}
-D\,\partial_rU=h_m\bigl(U_s-C_{\mathrm{env}}(t)\bigr),\qquad
-k\,\partial_rT=h\bigl(T_s-T_a(t)\bigr),
\label{eq:边界}
\end{equation}
$q=0$ 处取正则对称条件；初值 $U\equiv2.55$、$T\equiv301.15$ K。
式~\eqref{eq:质量坐标} 与边界式~\eqref{eq:边界} 的全部符号见表~\ref{tab:符号}；
四问的差别仅在取法（表~\ref{tab:设定}）：

\begin{table}[htbp]
\centering
\caption{问题 1--4 的取法对照}
\label{tab:设定}
\small
\begin{tabular}{lp{0.17\linewidth}p{0.16\linewidth}p{0.17\linewidth}p{0.22\linewidth}}
\toprule
取法 & 问题 1 & 问题 2 & 问题 3 & 问题 4 \\
\midrule
材料参数 & 附录 2（常数，$D=f(C)$） & 附录 3 & 附录 3 & 附录 4 \\
热湿耦合 & 解耦（$D$ 不含 $T$） & 双向耦合 & 双向耦合 & 双向耦合 \\
几何 & 固定 $R=0.02$ m & 固定 & 固定 & $R(t)$ 内生（半经验关系式，附件 2 交叉验证） \\
环境 & 附件 1（0--1800 s） & 附件 1 + 外推 & 同左 & 同左 \\
结果形式 & 场（表 1/2，result1） & 场（表 3/4，result2） & 决策量 $t_f$（表 5，result3） & 决策量 $t_f$（表 6，result4） \\
\bottomrule
\end{tabular}
\end{table}

\textbf{数值方法}：固定 $q$ 网格单元中心有限体积（N=160 主网格），内面通量共享；
表面用\textbf{半单元内阻与外对流阻串联}，表面值 $U_s,T_s$ 由
$8D_{\mathrm{eff}}(U_s-U_N)/(R^2\Delta q)=2h_m(C_{\mathrm{env}}-U_s)/R$ 联立重建，
其中 $D_{\mathrm{eff}}$ 取半单元中点值，不把末单元平均值当作表面值。
这样取值可避免强变的 $D\propto\mathrm e^{-0.45/C}$ 在端点取值时
产生内通量被强制为零的伪根。
\subsection{问题 1：预热平衡段的场（0--1800 s）}
\textbf{模型}：材料参数取附录 2 的常数值，热湿解耦，热方程为线性圆柱 Robin 问题；
湿分在 1800 s 内变化小但非零，仍按同一框架求解。
\textbf{求解与局部检验}：数值解与解析谱解逐点比较；后者取圆柱 Robin 特征展开
加时变环境的 Duhamel 卷积，并用稳态减法消除级数尾巴，
其自身误差 $<10^{-8}\,^{\circ}$C。单元中心最大偏差为 $5.9\times10^{-7}\,^{\circ}$C
（N=320），收敛阶为 2.00/1.99/1.99，见图~\ref{fig:解析验证}。
这一比较说明数值解确实求出了这组方程的准确解；由于两边用的是同一套材料参数，
它不能说明材料参数的取值是否正确。
\textbf{结果}：图~\ref{fig:温湿剖面} 给出温度与水分浓度的径向分布（色阶按时间递增），
图~\ref{fig:温度云图} 给出时空云图，图~\ref{fig:表面通量} 给出特征位置的时程与偏差、
表面通量与表面扩散系数。结果表见表~\ref{tab:表1}、表~\ref{tab:表2}（result1.xlsx）。

\begin{figure}[htbp]
\centering
\includefig[\linewidth]{fig12_温度场时空云图.pdf}{温度场时空云图}
\caption{温度场时空演化：云图以时间为横轴、径向位置为纵轴，
表面在约 600 s 内贴近环境、中心滞后明显，1800 s 末全场温差收敛到
约 $1\,^{\circ}$C 量级}
\label{fig:温度云图}
\end{figure}

\begin{figure}[htbp]
\centering
\includefig[\linewidth]{fig14_表面通量与表面扩散系数.pdf}{表面通量与表面扩散系数}
\caption{表面通量与表面扩散系数：(a) 表面水分
通量 $j_s$ 与表面扩散系数 $D_s=D(C_s)$ 的时程：通量初期最大而后衰减，
$D_s$ 随表面变干持续下降；(b) $D_s$ 与 $C_s$ 的对应关系，本过程的 $C_s$
只落在高含水端，故曲线平缓；图~\ref{fig:扩散系数}(a) 给出
$D\propto\mathrm e^{-0.89/C}$ 在全量程上的强非线性}
\label{fig:表面通量}
\end{figure}

\subsection{问题 2/3：全程场与固定几何达标时间}
\textbf{模型}：材料参数取附录 3，热与湿双向耦合，几何固定；达标判据
\begin{equation}
t_f=\inf\{t:\ \max_{q\in[0,1]}U(q,t)<0.15\ \text{kg/kg}\},
\label{eq:达标}
\end{equation}
用未舍入值在全域重建场上判定，判据即式~\eqref{eq:达标}。
本文不把中心最慢当作先验假设，而由逐步数值扫描确认：
全部时刻的最大值均位于中心列，判定容差取 $10^{-12}$。
\textbf{结果（场）}：温湿剖面随时间的演化（色阶按时间递增，viridis）见图~\ref{fig:温湿剖面}，
水分场时空云图与干壳/湿核结构见图~\ref{fig:水分云图}；干燥锋面自表面向内推进、
表面先干、中心最慢。随表面变干，$D\propto\mathrm e^{-0.45/C}$ 最多缩小 16.8 倍，
形成低扩散率“干壳”，中心残余水分只能穿过它排出。尾段 $\ln(U_{\max}-C_e)$
呈指数衰减（$R^2=0.9972$），衰减率对应剖面加权等效
$D\approx2.5\times10^{-10}$ m$^2$/s。

\begin{figure}[htbp]
\centering
\includefig[\linewidth]{fig3_温湿剖面.pdf}{温湿剖面}
\caption{温湿剖面演化}
\label{fig:温湿剖面}
\end{figure}

\begin{figure}[htbp]
\centering
\includefig[\linewidth]{fig13_水分场时空云图.pdf}{水分场时空云图}
\caption{水分场时空演化}
\label{fig:水分云图}
\end{figure}

\textbf{结果（达标时间）}：数值上报告\textbf{阈值穿越点所在的区间，右端加安全端点}：
问题 3 为 $[205837.7$, $205847.7]$ s，$\bm{t_f=57.1799}$ \textbf{h}，
插值值 57.1795 h，大于常系数参考时标 15.70 h，满足该判据。
收敛序列（N=40/80/160）为 57.249/57.201/57.180 h，
Richardson 外推 57.1636 h，主值偏差 0.016 h 属空间误差量级。
四位小数舍入后的 60 s 网格判定为 57.2333 h，与未舍入值差 0.053 h，
属另一种数据使用方式的差异，见第七章。阈值 0.15 干基 $=13.04\%$ 湿基，
与药材水分限度的通行标准（约 13\%）相符~\cite{pharmacopoeia}。
结果表见表~\ref{tab:表3}、表~\ref{tab:表4}（result2.xlsx）与
表~\ref{tab:表5}（result3.xlsx）；达标时刻见图~\ref{fig:穿越}。

\textbf{潜热对照情景}：第三章假设 3 的无潜热近似用情景检验界定：表面相变汇
$-k\partial_rT|_R=h(T_s-T_a)+L_vj_w$，
$j_w=\rho_{d,s}h_m(U_s-U_e)$（$\rho_{d,s}=(650+128U_s)/(1+U_s)$ 为表面局部
干物质密度，$L_v=2.4\times10^6$ J/kg 近似值）。对照情景沿用同一套附录 3 材料参数、
同一边界数据与同一数值框架，$t_f$ 由 57.1799 h 变为 \textbf{62.5694 h}
（$\bm{+5.3889}$ \textbf{h}，$+9.4\%$），全程最大逐时表面/中心温差
$29.6/29.2\,^{\circ}$C、最大平均含水率差 $0.322$ kg/kg。故无潜热基准的 $t_f$ 偏差量级约 5.4 h。
该情景中表面温度一度降至 $13.4\,^{\circ}$C、低于湿球温度，
说明 $h_m$ 与能量边界条件尚需协调。

\subsection{问题 4：收缩移动域与内生收缩模型}
\textbf{几何与移动域}：仿射收缩 $r=R(t)\sqrt q$ 下方程同形，见式~\eqref{eq:质量坐标}，
$R(t)$ 为待确定的几何函数。\textbf{附件 2 的三段结构}是 A 急缩 0--3.5 h 占 61\%、
B 缓缩 3.5--21 h、C 近平台段约 21 h 起，它不是单一失水收缩曲线：
理想等容收缩这一特定基准 $R_0\sqrt{(1+\bar U)/(1+C_0)}$ 不能解释全部数据，
其反演平均含水率为 0.278 kg/kg，低于模型同刻的 0.149 kg/kg，
在 $0.129\to0.155$ kg/kg 区段还存在成孔造成的偏差，并与数据在约 16 h 交叉滞后，
附件 2 的时间范围是 0--72 h，$t_f$ 落在该区间内，见图~\ref{fig:收缩机制}。

\begin{figure}[htbp]
\centering
\includefig[\linewidth]{fig6_收缩机制分析.pdf}{收缩机制分析}
\caption{附件 2 的收缩机制分析。附件 2 呈 A 急缩 61\%、B 缓缩、C 平台的三段复合结构，
与约 16 h 滞后交叉；不作参数拟合的“理想收缩+玻璃化冻结”半经验关系式
已能由结构给出减速与平台，其中按平均含水率判断结皮的那一条给出平台 1.169 cm、
偏差 $-2.4\%$、冻结时刻偏差 $+12\%$；
早期塌陷与后期成孔对应两组机制。数据源：问题 4 主解湿度场与附件 2 实测半径}
\label{fig:收缩机制}
\end{figure}

\textbf{半经验收缩关系式}：问题 4 主模型由三项物理机制叠加而成，
每项对应一种物理过程，半径由模型自身逐步预测，共 5 个参数：
\begin{equation}
R(t)=\underbrace{R_0\sqrt{\frac{1+\bar U(t)}{1+C_0}}}_{\text{失水收缩}}\cdot
\underbrace{\bigl[1-\varepsilon_c\,g_c(\bar U)\bigr]}_{\text{毛细-膨压塌陷}}\cdot
\underbrace{\bigl[1+\varepsilon_p\,g_p(C_s)\bigr]}_{\text{结皮-成孔停滞}}
\label{eq:收缩闭合}
\end{equation}
(i) 失水收缩：这一项不含待定参数，由水分守恒与固体、水等密度的假设给出；
(ii) 毛细-膨压塌陷：$g_c(x)=(1-x)^{a}x^{b}$（$x=(C_0-\bar U)/C_0$，峰值归一），
其量级由毛细压实的估计支撑，估计区间 $[0.028,0.580]$ cm
包含了 2 h 观测到的 0.182 cm 超额与 9.1\% 的半径应变；
(iii) 结皮-成孔停滞：表面 $C_s$ 达到 $C_{\mathrm{glass}}\approx0.21$ 后收缩停滞，
继续失水转化为孔隙，这一阈值取 Gordon--Taylor 结构的典型参数情景值~\cite{saravacos}，
已列入参数不确定清单；
$g_p=\mathrm{clip}\bigl((C_{\mathrm{glass}}-C_s)/(C_{\mathrm{glass}}-C_{s,\min}),0,1\bigr)^{k_p}$，
$C_{s,\min}=0.0488$ kg/kg 为完整轨迹读得的归一化常数。
$g_c,g_p$ 是有物理含义的经验形状函数，该式整体为\textbf{半经验}关系式；
5 个参数由最小二乘对附件 2 标定，$C_{\mathrm{glass}}$ 与 $\rho_{sk}=1468$ 这 2 个量
取情景值或典型范围值，不参与拟合，参数表见附录 B，标定细节见附录 C。

\textbf{标定与验证}：拟合优度 RMSE $=0.047$ mm，约为附件 2 量化分辨率 0.01 mm 的
5 倍，最大偏差 0.162 mm；滞后交叉约 17 h（计算值 17.0 h），数据为约 16 h；
平台高度 1.1969 cm，数据为 1.198 cm；三段形态全部复现，见图~\ref{fig:内生贴合}。
孔隙份额自 9.5 h 起转为正值并单调增加，成孔由此开始，末期孔隙率为 30.9\%。
过拟合检验表明：参数标定与验证用的是同一批数据，
$R_{\mathrm{obs}}\to(\bar U,C_s)_{\text{主解}}\to\text{标定}\to R_{\mathrm{pred}}$
构成循环；奇偶折半留出法只检验了同一条光滑曲线上的\textbf{插值稳健}；
若要独立验证，需对每组候选参数从相同初值重解，即\textbf{完整耦合标定}，
本文当前的做法是借助实测几何生成湿度场再作标定，尚属近似。

\begin{figure}[htbp]
\centering
\includefig[\linewidth]{fig7_内生收缩贴合.pdf}{内生收缩贴合}
\caption{内生收缩模型对附件 2 的吻合与验证：式~\eqref{eq:收缩闭合} 由三项物理机制
叠加而成，含 5 个标定参数与 2 个不参与拟合的量，它复现了附件 2 的全程形态：
RMSE 0.047 mm，约为量化分辨率的 5 倍；滞后交叉 17.0 h，数据约 16 h；
平台 1.1969 cm，数据 1.198 cm；三段形态全部复现。
残差子图显示偏差始终在 $\pm0.16$ mm 以内，这是以实测几何单向驱动得到的拟合指标}
\label{fig:内生贴合}
\end{figure}

\textbf{耦合主解}：附件 2 在此仅作验证数据。主结果采用完全内生的收缩，
$R$ 由湿度场逐步预测，每步以当前 $\bar U,C_s$ 代入式~\eqref{eq:收缩闭合}，
不使用附件 2 的几何数据，得 $\bm{t_f=50.5369}$ \textbf{h}，该值取自 result4 内生版的
第 3033 行；未对右端点舍入，达标时刻所在区间为 $[181917.7,181932.7]$ s，
瞬时首模近似给出的参考值为 19.02 h，该结果满足这一参考。
把附件 2 实测几何作为输入的对照版得 $t_f=50.8230$ h，与主解的偏差为
$-0.286$ h、$-0.56\%$，这一比较属模型之间的对照，不是独立的实验验证。
两个精度指标须分清：\textbf{0.047 mm 是以实测几何单向驱动得到的拟合指标}，
\textbf{耦合主解自身的指标}是沿满 72 h 真实轨迹的全程 RMSE \textbf{0.088 mm}，
其中 0--达标 \textbf{0.065}、达标--72 h \textbf{0.127} mm；
主解积分满 72 h，首次达标单独记录，未用常值填充。
末期 $R_{\mathrm{pred}}(72\,\mathrm h)=1.1836$，附件 2 为 1.198 cm，
60--72 h 再降 0.031 mm，末端已近似平台，但仍有残余缓降，与结构上的限定一致。
半径预测所用代数关系的残差为：先在 $R_{\mathrm{ideal}}$ 下求出 $C_s$，
再代入修正式后未收敛的残差，其最大值 0.003 mm；末单元温度与重建表面温度
之差最大 0.018 ${}^\circ$C，对 $t_f$ 的影响远在 0.01 h 以下。
收敛序列 50.5450/50.5370/50.5369 h 的相邻差为 28.8 s 与 0.41 s，
后者已小于尾段约 15 s 的步宽，故不能用表观收敛阶判断收敛；
改用连续事件定位的 $t^*$ 序列为 50.5422/50.5355/50.5362 h，
N 由 80 增至 160 时变化 0.0007 h，可视为空间误差的量级，见第七章。
中心/表面演化与半径收缩见图~\ref{fig:中心表面}，截面演化见图~\ref{fig:截面}；
结果表见表~\ref{tab:表6}（result4.xlsx）。

\begin{figure}[htbp]
\centering
\includefig[0.86\linewidth]{fig16_中心表面演化与半径收缩.pdf}{中心/表面演化与半径收缩}
\caption{中心/表面含水率演化与半径收缩：(a) 中心与表面含水率曲线，
配 0.15 kg/kg 阈值虚线与 $t_f=50.5369$ h 达标竖线：表面先到干态、中心最后达标；
(b) 内生半径 $R(t)$ 在 50.5 h 进入近平台，60--72 h 仅再降 0.031 mm，
模型已积分至 72 h，曲线止于 $t_f=50.54$ h}
\label{fig:中心表面}
\end{figure}

\begin{figure}[htbp]
\centering
\includefig[\linewidth]{fig9_截面演化.pdf}{截面演化}
\caption{问题 4 的截面演化。图中给出 6 个代表时刻的圆盘截面热图，
色标为水分浓度，取对数色阶以便分辨低含水端的壳核差异，
色标上 0.15 处即题目的达标阈值。热图显示干燥自表面向中心推进，
外半径同步收缩并进入平台：附件 2 实测的平台半径为 1.198 cm，
模型逐步预测的末端半径为 1.19 cm，各面板标题给出该时刻的模型半径。
图底的时间轴带给出内生半径 $R(t)$ 与六个采样时刻在曲线上的位置，
六个盘的时间间隔是 5/10/10/15/10 h，而半径降幅集中在前两盘：
$R$ 在前 5 h 由 2.00 降到 1.42 cm，占 0.81 cm 总降幅的 72\%，
20 h 后即进入平台，因此在等间隔取样、等宽排版下视觉变化前快后慢；
中心为最慢点，$t_f=50.5369$ h 时中心首次低于 0.15 kg/kg}
\label{fig:截面}
\end{figure}

\textbf{与问题 3 的大小关系}：
\begin{itemize}
  \item \textbf{附录材料参数组合效应}：附录 4 与附录 3 的 $D$ 之比
        $0.175\,\mathrm e^{0.15/U}$ 由 $1/5.39$（$U=2.55$）变到 $1/2.10$（$U=0.15$）；
        固定半径下 $t_f$ 从 57.18 h 放大到 129.10 h，即 $\times2.26$、$+78.57$ h。
        该对照同时改变了全部热物理参数，因此不能严格归因于单个扩散系数的倍率。
  \item \textbf{收缩效应}：$R^2$ 降 64\%，扩散时间尺度缩至 0.36；$2/R$ 换热与传质
        面积比增 67\%，把 129.10 h 拉回 50.54 h，即 $-60.9\%$。
        三组对照中，收缩带来的加速幅度最大。
  \item \textbf{净效应}：材料参数组合效应带来的减速 $\times2.26$ 被收缩加速反超，
        问题 4 反而比问题 3 快 $11.6\%$，比固定半径的对照快 $+6.64$ h。
\end{itemize}

\subsection{力学驱动扩展：收缩的力学机制}
本节把半经验关系式推广为“由守恒律与有限应变力学约束、材料本构可标定的
内生收缩模型”，作为\textbf{扩展研究}；主线为 5.4 节的半经验关系式，
它对附件 2 的精度更高；

\textbf{框架}：材料点 $X$ 的现时位置 $r(X,t)$ 为未知量，有限应变非仿射几何
$\lambda_r=r_X$、$\lambda_\theta=r/X$、$\lambda_z=1$、$J=\lambda_r\lambda_\theta$；
孔隙非负直接进入方程：$g=J-s_0-w\ge0$ 由乘子 $\zeta$ 以互补条件活动集实施；
毛细储能 $\psi_{\mathrm{cap}}=(J-s_0)\int_S^1p_c\,ds$，保水本构取幂律
$p_c=p^*(1-S)/S^\beta$，$\bar p_c$ 有解析式，与数值积分的偏差
不超过 $5\times10^{-16}$；骨架取对数应变的单松弛黏弹性模型，
含内变量 $Z$，用后向欧拉作局部消元，每个时间步只需解一次；
初始平衡 $e_*=\bar p_c(S_0)/K_\infty$ 使初始总应力为零；
总应力 $\sigma_i=H_i/J+\bar p_c-\zeta$，准静态径向平衡以弱式求解，
表面取 $\sigma_{rr}=0$，\textbf{外半径 $R(t)$ 是待求量而非已知输入}；
与传质以 $U_t=\partial_q(D/r_q^2U_q)$ 衔接，
每步按力学、几何系数、传质的顺序迭代耦合；热算子按
$a_{\mathrm{eff}}T_t=(1/(rr_q))\partial_q(rk/r_qT_q)$ 的有限体积形式离散。
判别实验表明应力传递是全局的，而不是逐点体积平衡：
只干燥最外层时，湿内部 $J$ 保持 1.000，而径向应力以 $-0.047$ 全域均匀传递，
表皮环向应力为 $+4.188$，表面节点内移 0.042 cm。

\textbf{退化检查}（表~\ref{tab:力学退化}）：孔隙非负全程成立，
体积衡等式达到机器精度，热算子与参考解的比较结果一致。
在常数材料参数、无失水、阶跃环境下取同一后向欧拉步，N=16/80/160 的最大场差为
$0.0084/0.0023/0.0013\,^{\circ}$C，并随加密单调下降；牛顿步全部收敛，
最大平衡残差为 $9.5\times10^{-11}$。

\begin{table}[htbp]
\centering
\caption{力学模型的退化检验（末列为实测判定）}
\label{tab:力学退化}
\small
\begin{tabular}{p{0.26\linewidth}p{0.26\linewidth}p{0.31\linewidth}l}
\toprule
检查 & 应出现的结果 & 实测 & 判定 \\
\midrule
不干燥 + 初始总应力平衡 & 尺寸无人为跳变 & $\max|\Delta r|=0.000$ & $\checkmark$ \\
不计黏弹分支 & 恢复纯弹性问题 & 计算结果正常（$R_{\mathrm{end}}=1.717$ cm） & $\checkmark$ \\
$\tau\ll$ 过程尺度 & 黏弹分支充分松弛 & 与纯弹性差 $5.3\times10^{-15}$ cm & $\checkmark$ \\
$\tau\gg$ 过程尺度 & 内变量近冻结 & $\max|Z|/\max|E|=4.3\times10^{-11}$ & $\checkmark$ \\
指定仿射映射 $r=R\sqrt q$ & 传质系数恢复 $4qD/R^2$ & 相对误差 $4.9\times10^{-15}$ & $\checkmark$ \\
固定外形持续失水 & $g$ 按体积恒等式增长 & 相对误差 $2.2\times10^{-16}$ & $\checkmark$ \\
全程孔隙非负 & 任一时刻任一点 $g\ge0$ & $g_{\min}=+0.1446$（无任何负孔隙） & $\checkmark$ \\
全局体积核验式 & $R^2/R_0^2=s_0+\frac{\rho_{d0}}{\rho_w}\bar U+\int g\,dq$ & 残差 $\le2.7\times10^{-16}$ & $\checkmark$ \\
热算子退化比对 & 与参考解同 BE 步一致，随加密下降 & $0.0084/0.0023/0.0013\,^{\circ}$C & $\checkmark$ \\
\bottomrule
\end{tabular}
\end{table}

\textbf{标定与复现程度}：幂律保水—骨架结构对附件 2 标定，
$\beta$ 的下界为 0.02，4 个起点都收敛到同一解：
$p^*=1.0\times10^{4}$ 触到上界，此时轨迹由无量纲形状
$\bar p_c(S)/\bar p_c(S_0)$ 决定，$p^*$ 处于渐近不变区，不可辨识属于预期；
$\beta=0.1059$；$\Pi_\tau=1.0\times10^{-3}$ 触到下界。
完整积分满 72 h 的 \textbf{RMSE $=0.522$ mm}，其中急缩段 1.193、
缓缩段 0.875、平台段 0.163 mm；分段为 0--达标 0.604、达标--72 h 0.212 mm。
末期平台 1.1747，附件 2 为 1.198 cm，60--72 h 仅再降 0.04 mm，
末端已真正进入平台；$\bm{t_f=51.1667}$ \textbf{h}，
较半经验主线 50.5369 h 偏差 $+1.25\%$。
奇偶折半留出为 0.587/0.591 mm，说明没有过拟合；分块自助法给出 $\beta$ 的 95\%
置信区间 $[0.095,\,0.91]$；孔隙非负全程成立，$g_{\min}=0.1446$。
平台来自\textbf{吸力饱和平衡}：标定参数下表面 Deborah 数
De $\approx1.0\times10^{-3}\ll1$，说明平台并非黏弹冻结所致，见图~\ref{fig:力学升级}。

\begin{figure}[htbp]
\centering
\includefig[\linewidth]{fig10_力学升级.pdf}{力学升级}
\caption{力学驱动内生收缩模型：修正后的数值格式与幂律保水重标定（$\beta=0.106$）
对附件 2 的拟合结果为 RMSE 0.522 mm，残差在三段上分别为急缩 1.193、
缓缩 0.875、平台 0.163 mm，$t_f=51.1667$ h；末端 60--72 h 仅再降 0.04 mm，
已真正进入平台，平台机制为吸力饱和平衡，此时 De $\approx10^{-3}\ll1$。
(c) 面板给出该平台机制：归一化平均吸力 $\bar p_c/p^*$ 随饱和度 $S$ 下降而单调上升，
并在 $S\to0$ 时趋于有限上限 $p^*/[(1-\beta)(2-\beta)]$，
末端收缩进入平台即来自这一吸力饱和平衡；
气相体积份额全程非负，$g_{\min}=0.1446$，见表~\ref{tab:力学退化}}
\label{fig:力学升级}
\end{figure}

\textbf{由模型结构决定的上限与参数缺失}：在所选幂律保水—简化储能结构内，
$\bar p_c$ 的有限上限为 $p^*/[(1-\beta)(2-\beta)]$，$\beta=0.15$ 时仅为 $0.636\,p^*$；
且吸力的作用范围是局部的，只作用于干燥表皮，
湿内部则以其体积刚度抵抗由外层传递而来的压缩。由此带来两个后果：
第一，附件 2 前 21 h 的早期全局塌陷无法仅用吸力驱动的局部收缩解释；
第二，末端平衡 $J_{\mathrm{end}}(\beta)$ 由 $J\cdot g(S(J);\beta)=g(S_0;\beta)$ 决定，
$\beta=0.02/0.05/0.10/0.15/0.85$ 时 $J_{\mathrm{end}}=0.368/0.356/0.336/0.319/0.218$。
若要达到附件 2 的平台 $J=0.36$，须取 $\beta\lesssim0.05$，此时缓缩段过慢；
若要让中段足够快，须取 $\beta\gtrsim0.10$，此时末端必然过缩。
因此在这一结构内，单一 $\beta$ 无法同时给出快速的急缩与停滞的平台；
标定折中取 $\beta=0.106$，末端 $J=0.3343$，附件 2 的对应值为 0.36。
这一上限由所选本构与储能结构决定，其他保水曲线或体积驱动机制有可能达到
更高的精度，这一问题留作后续工作。\textbf{参数缺失}方面，真实保水曲线
$p_c(S,T)$ 是其中影响最大的一项，此外还有膨压类体积驱动、干态模量增长与结皮硬化、
骨架模量与松弛谱、初始湿质量与真实密度；体积核算只给出 $\rho_{d0}\le256$ kg/m$^3$
的上限。

\textbf{结论}：当前的精度仍不如半经验主线，其指标为 0.047/0.088 mm。
在真实保水曲线与力学实测数据到位之前，本文的力学模型还不能复现附件 2 的实测过程；
力学驱动的推广是后续可继续沿用的方向，届时半经验收缩关系式可作为其低维验证基准。

\section{模型检验}

\subsection{无量纲判据（$Le$、$Bi_m$、$Bi_h$）}
沿矩形扫描域 $C\in[0.15,2.55]$、$T\in\{301.15,\,323.315\}$ K、$R=R_0$，逐点
按附录公式计算 $\alpha=k/(\rho c_p)$，扫描区间即表~\ref{tab:判据族} 所列区间。
干燥末段表面 $U<0.15$ 的局部状态超出该扫描域，扫描未包含。
表中回代偏差指采用该简化后 $t_f$ 与全耦合解 57.2009 h 的实测差。

\begin{table}[htbp]
\centering
\caption{无量纲判据表：定义、扫描区间、阈值语义、简化决策与回代偏差}
\label{tab:判据族}
\small
\begin{tabular}{llp{0.16\linewidth}p{0.16\linewidth}p{0.23\linewidth}l}
\toprule
判据 & 定义 & 附录 3 区间 & 附录 4 区间 & 简化决策 & 回代偏差 \\
\midrule
$Le$ & $\alpha/D$ & $10.6860$--$640.6171$ & $31.2068$--$556.1413$ & $Le\gg1$：热场准稳态 $T\approx T_a$ & $0.303$ h \\
$Bi_m$ & $h_mR/D$ & $1.1805$--$47.7480$ & $6.3606$--$100.3742$ & $Bi_m\to0$：集总（内扩散无阻力） & $-46.03$ h \\
$Bi_h$ & $hR/k$ & $1.0353$--$1.9263$ & $1.8964$--$3.4226$ & $Bi_h\to0$：忽略表面热滞后 & $0.075$ h \\
\bottomrule
\end{tabular}
\end{table}

回代检验结果如下，每条给出采用该简化后 $t_f$ 的偏移：
\begin{itemize}
  \item \textbf{$Le\gg1$ 说明热场先趋于准稳态}：温度场取 $T_a(t)$ 只求湿分，
        $t_f$ 偏差 $0.303$ h（0.5\%）；本文主模型仍保留双向耦合，此简化只作
        回代检验，且只在无潜热基准内评价，不适用于有蒸发冷却的情形
        （5.3 节潜热情景表面温差达 $29.6\,^{\circ}$C）。
  \item \textbf{$Bi_m$ 区间 $\gtrsim O(1)$ 表明内部扩散主导}：集总模型给出
        $t_f=11.17$ h，比全耦合解快 $46.03$ h（$-80\%$）；可见把含水率作
        “集总或平均”处理都无法复现本题的 $t_f$，模型必须保留空间分布。
  \item \textbf{$Bi_h\sim O(1)$ 下达标时间对该简化不敏感}：令 $h\to\infty$
        实测偏差 $0.075$ h（0.13\%），本文未采用该简化；$Bi_h\sim1$ 本身
        并不表示空间温差小，同样不适用于有蒸发冷却的情形。
\end{itemize}

\subsection{参考时标（常系数与瞬时首模近似）}
取 $D$ 的路径点态最大 $D_{\max}$（$D$ 对 $C,T$ 均单调增），构造常系数 Robin
参考问题，其首模形式为
\begin{equation}
t_{\mathrm{ref}}=\frac{R^2}{D_{\max}\lambda_1^2}\,
\ln\!\frac{A_1\,(C_0-C_e)}{C_{\mathrm{th}}-C_e},
\label{eq:时标}
\end{equation}
其中 $\lambda_1$ 为 $\lambda J_1(\lambda)=Bi_m J_0(\lambda)$ 第一根（圆柱 Robin
特征值问题，标准谱理论见~\cite{crank}），$A_1$ 为中心响应系数。
问题 3（式~\eqref{eq:时标}）：$D_{\max}=1.359\times10^{-8}$ m$^2$/s、$Bi_m=1.177$，得
$t_{\mathrm{ref}}=15.70$ h（300 模精确级数的达标时刻同为 15.7021 h，参考问题
自身的首模近似误差可忽略）；膜界 $11.17$ h（忽略内部扩散阻力，系统性偏弱，
仅作次级参照）。问题 4：把瞬时首模时标沿真实 $R(s)$ 轨迹积分，得参考值
$19.02$ h（$R_{\min}$ 常值版 15.95 h、固定 $R_0$ 版 37.61 h）。
两个主值（57.18 h、50.54 h）均远大于各自参考时标。

\textbf{参考解与真实问题之间的差异}：(i) 参考解把环境平衡势固定为
末值 $C_e=0.04986$，但真实边界初期低至 0.01963——同一 $D_{\max}$ 首模参考
在固定初值边界下为 \textbf{14.5491} h，用于排除简化的时标随边界取法而变，
比较的先后顺序并不固定；(ii) 对散度形式的变系数扩散方程，
$D\le D_{\max}$ 本身不构成逐点解比较；(iii) 收缩版积分瞬时第一特征值，
而 $Bi(s)$ 随时间变化引起模态耦合，瞬时特征值积分不是精确的变边界谱解。
若达标时间远低于参考时标，就应怀疑该结果是否可信。参考时标只能用于判断量级，
据此可以排除偏差较大的简化。

\subsection{检验结果}
解析谱解与守恒有限体积数值解来自两种不同的数学结构，
两者的逐点最大偏差为 $5.9\times10^{-7}\,^{\circ}$C（V1）。
各项检验汇总于表~\ref{tab:验证}，
收敛与守恒的可视化见图~\ref{fig:收敛守恒}。


\begin{table}[H]
\centering
\caption{验证汇总表（V1--V6）：手段、数学结构、实测数字、能得到的结论与尚未检验的误差源}
\label{tab:验证}
\footnotesize
\begin{tabular}{p{0.30\linewidth}p{0.30\linewidth}p{0.32\linewidth}}
\toprule
检验（手段） & 实测数字 & 能得到的结论 / 尚未检验的误差源 \\
\midrule
V1 问题 1 热场解析谱解（特征展开 + Duhamel 卷积）与有限体积数值解 &
单元中心逐点最大偏差 $5.9\times10^{-7}\,^{\circ}$C（N=320）；收敛阶 2.00/1.99/1.99 &
数值解求的正是该线性模型，网格加密后按二阶收敛，场值不是网格假象；材料参数的取值不在本文检验范围内 \\
V2 后验尾段诊断（等温线性化 Bessel 首模结构对照） &
尾段 $\ln(U_{\max}-C_e)$ 线性 $R^2=0.9972$；衰减率 $\gamma=3.55\times10^{-6}$/s 对应
等效 $D=2.54\times10^{-10}$（$C\approx0.108$，在尾段剖面区间内） &
尾段的衰减确实由扩散首模控制（等效 $D_{\mathrm{eff}}$ 由本解反演，属构造一致、
不是独立预测）；尚未检验的边界与几何离散误差 \\
V3 参考时标判据检验（常系数/瞬时首模近似，见 6.2 节限定） &
$t_f=57.18$ h $>$ 15.70 h；$t_f=50.54$ h $>$ 19.02 h &
达标时间与参数给出的参考量级一致（远大于参考时标），据此可以排除偏差较大的简化；
尚未检验的变系数扩散与变动边界谱差异 \\
V4 守恒恒等式（总水量、冻结热容的显热离散平衡，逐项核对） &
水量窗口相对残差 $\le5.5\times10^{-11}$；显热离散平衡全局吞吐残差 $\le1.1\times10^{-11}$ &
逐项核对说明边界输入的水量与内部贮存的变化一致，水分与显热不会凭空增减，
达标时间不是数值泄漏造成的；状态焓未独立定义，
真实能量守恒与潜热分析不在本文检验范围内 \\
V5 退化测试（移动域 $R(t)\equiv R_0$；$D$ 冻结） &
移动域退化为固定域：$\Delta t_f=0.00$ s、场差 0；$D$ 冻结退回 Bessel 首模：偏差
$3.1\times10^{-9}$、首模商 0.99996 &
移动域在 $R(t)\equiv R_0$ 时还原为固定域，说明两个模型自洽；常系数湿场退回
Bessel 首模，是独立对照。模型的基本假设本身不在本文检验范围内 \\
V6 网格/时间 Richardson 收敛 &
场级空间阶 1.90--2.45；$t_f$ 序列 57.249/57.201/57.180 h，外推偏差 0.016 h；
时间容差收紧十倍实测 $\Delta t_f=0.0014$ h &
报告的 $t_f$ 是连续方程的解，不是网格假象；单独使用不足以支撑可靠性 \\
\bottomrule
\end{tabular}
\end{table}

\begin{figure}[htbp]
\centering
\includefig[0.86\linewidth]{fig2_问题1解析验证.pdf}{问题 1 解析验证}
\caption{问题 1 解析比较（V1）：(a) 中心与表面两个特征位置的温度数值解
与解析谱解时程；(b) 两个特征位置的逐点偏差时程，量级
$10^{-6}\,^{\circ}$C、全程低于 $3.2\times10^{-6}\,^{\circ}$C，虚线为四位小数输出的
分辨率上限；径向网格加密的二阶收敛见表~\ref{tab:验证} 与
图~\ref{fig:收敛守恒}(a)；全域单元中心在 N=320 的最大偏差为
$5.9\times10^{-7}\,^{\circ}$C}
\label{fig:解析验证}
\end{figure}

\begin{figure}[htbp]
\centering
\includefig[\linewidth]{fig15_收敛与守恒检验.pdf}{收敛与守恒检验}
\caption{收敛与守恒检验：(a) 温度场误差随网格加密按二阶下降（V1、V6）；
(b) $t_f$ 随网格加密的收敛序列；
(c) 残差随网格加密：水量与显热离散平衡残差维持在 $10^{-11}$ 量级附近（V4）}
\label{fig:收敛守恒}
\end{figure}

\section{结果分析与敏感性}

结果表格按 result1--4.xlsx 汇总，见表~\ref{tab:表1}--表~\ref{tab:表6}；
达标时刻的三条曲线分别取问题 3 主解、问题 4 主解与附录 4 固定半径对照，
见图~\ref{fig:穿越}。

% 本文件由 gen_tables.py 从 03-数据/*.csv 自动生成，禁止手工改数。

\begin{table}[htbp]
\centering
\caption{问题~1（预热平衡阶段）药材温度（℃）。数字由 result 工作簿同一数据源抽取，保留四位小数。}
\label{tab:表1}
\small
\begin{tabular}{lccccc}
\toprule
时间/s & 0 & 0.5 & 1 & 1.5 & 2 \\
\midrule
100s & 28.0001 & 28.0003 & 28.0040 & 28.0327 & 28.1801 \\
300s & 28.0408 & 28.0635 & 28.1514 & 28.3681 & 28.8489 \\
600s & 28.4533 & 28.5360 & 28.8039 & 29.3159 & 30.1652 \\
900s & 29.3243 & 29.4583 & 29.8755 & 30.6161 & 31.7304 \\
1200s & 30.5427 & 30.7098 & 31.2223 & 32.1126 & 33.4276 \\
1500s & 31.9957 & 32.1867 & 32.7660 & 33.7463 & 35.1203 \\
1800s & 33.5753 & 33.7720 & 34.3642 & 35.3621 & 36.7856 \\
\bottomrule
\end{tabular}
\end{table}

\begin{table}[htbp]
\centering
\caption{问题~1（预热平衡阶段）药材水分浓度（kg/kg）。数字由 result 工作簿同一数据源抽取，保留四位小数。}
\label{tab:表2}
\small
\begin{tabular}{lccccc}
\toprule
时间/s & 0 & 0.5 & 1 & 1.5 & 2 \\
\midrule
100s & 2.5500 & 2.5500 & 2.5500 & 2.5500 & 2.2470 \\
300s & 2.5500 & 2.5500 & 2.5500 & 2.5492 & 2.0508 \\
600s & 2.5500 & 2.5500 & 2.5500 & 2.5353 & 1.8770 \\
900s & 2.5500 & 2.5500 & 2.5497 & 2.5045 & 1.7547 \\
1200s & 2.5500 & 2.5500 & 2.5482 & 2.4646 & 1.6586 \\
1500s & 2.5500 & 2.5499 & 2.5445 & 2.4206 & 1.5787 \\
1800s & 2.5500 & 2.5497 & 2.5383 & 2.3755 & 1.5102 \\
\bottomrule
\end{tabular}
\end{table}

\begin{table}[htbp]
\centering
\caption{问题~2（全烘干过程）药材温度（℃）。数字由 result 工作簿同一数据源抽取，保留四位小数。}
\label{tab:表3}
\small
\begin{tabular}{lccccc}
\toprule
时间/h & 0 & 0.5 & 1 & 1.5 & 2 \\
\midrule
0.5h & 32.1892 & 32.3821 & 32.9659 & 33.9608 & 35.4130 \\
1.0h & 40.3816 & 40.5536 & 41.0605 & 41.8782 & 42.9976 \\
1.5h & 45.8468 & 45.9348 & 46.1930 & 46.6051 & 47.1401 \\
2.0h & 48.4502 & 48.4881 & 48.5984 & 48.7735 & 49.0033 \\
2.5h & 49.4670 & 49.4792 & 49.5136 & 49.5653 & 49.6609 \\
3.0h & 49.8495 & 49.8553 & 49.8746 & 49.9101 & 49.9664 \\
\bottomrule
\end{tabular}
\end{table}

\begin{table}[htbp]
\centering
\caption{问题~2（全烘干过程）药材水分浓度（kg/kg）。数字由 result 工作簿同一数据源抽取，保留四位小数。}
\label{tab:表4}
\small
\begin{tabular}{lccccc}
\toprule
时间/h & 0 & 0.5 & 1 & 1.5 & 2 \\
\midrule
0.5h & 2.5499 & 2.5489 & 2.5256 & 2.3257 & 1.6486 \\
1.0h & 2.5258 & 2.4948 & 2.3578 & 2.0230 & 1.4711 \\
1.5h & 2.3861 & 2.3256 & 2.1344 & 1.8020 & 1.3476 \\
2.0h & 2.1709 & 2.1084 & 1.9236 & 1.6259 & 1.2311 \\
2.5h & 1.9566 & 1.9006 & 1.7360 & 1.4720 & 1.1166 \\
3.0h & 1.7662 & 1.7165 & 1.5702 & 1.3333 & 1.0081 \\
\bottomrule
\end{tabular}
\end{table}

\begin{table}[htbp]
\centering
\caption{问题~3 药材烘干过程的水分浓度（kg/kg）。结束行时刻为未舍入判定值 57.1799~h；表中 0.1500 为舍入显示，判定位用未舍入值。}
\label{tab:表5}
\small
\begin{tabular}{lccccc}
\toprule
时间/h & 0 & 0.5 & 1 & 1.5 & 2 \\
\midrule
6h & 1.0156 & 0.9868 & 0.9006 & 0.7546 & 0.5335 \\
12h & 0.4548 & 0.4418 & 0.4019 & 0.3289 & 0.1639 \\
18h & 0.2979 & 0.2906 & 0.2679 & 0.2247 & 0.0842 \\
24h & 0.2371 & 0.2321 & 0.2161 & 0.1851 & 0.0663 \\
30h & 0.2053 & 0.2013 & 0.1887 & 0.1637 & 0.0597 \\
36h & 0.1854 & 0.1821 & 0.1714 & 0.1501 & 0.0565 \\
42h & 0.1717 & 0.1687 & 0.1594 & 0.1404 & 0.0547 \\
48h & 0.1615 & 0.1588 & 0.1504 & 0.1332 & 0.0536 \\
54h & 0.1535 & 0.1511 & 0.1434 & 0.1275 & 0.0528 \\
烘干结束时间 & 0.1500 & 0.1477 & 0.1402 & 0.1249 & 0.0525 \\
\bottomrule
\end{tabular}
\end{table}

\begin{table}[htbp]
\centering
\caption{问题~4 药材烘干过程的水分浓度（kg/kg）。“—”表示该时刻该固定位置已在药材表面之外（半径收缩，留空不外推）。结束行时刻为未舍入判定值 50.8230~h。}
\label{tab:表6}
\small
\begin{tabular}{lccccc}
\toprule
时间/h & 0 & 0.5 & 1 & 1.5 & 药材表面 \\
\midrule
6h & 1.7164 & 1.5350 & 1.0212 & — & 0.4215 \\
12h & 0.7340 & 0.6516 & 0.4068 & — & 0.1666 \\
18h & 0.4065 & 0.3667 & 0.2387 & — & 0.0888 \\
24h & 0.2837 & 0.2598 & 0.1783 & — & 0.0671 \\
30h & 0.2253 & 0.2085 & 0.1488 & — & 0.0593 \\
36h & 0.1920 & 0.1790 & 0.1312 & — & 0.0558 \\
42h & 0.1706 & 0.1598 & 0.1196 & — & 0.0539 \\
48h & 0.1556 & 0.1463 & 0.1113 & — & 0.0529 \\
烘干结束时间 & 0.1500 & 0.1413 & 0.1081 & — & 0.0525 \\
\bottomrule
\end{tabular}
\end{table}

\begin{figure}[htbp]
\centering
\includefig[0.86\linewidth]{fig4_达标穿越.pdf}{达标时刻}
\caption{达标时刻：$\max_qU(q,t)$ 全程曲线与
0.15 阈值线；问题 3 主解 $t_f=57.1799$ h、问题 4 主解 $t_f=50.5369$ h（内生几何）
与附录 4 固定半径对照 129.1047 h 三条曲线对比，横轴标出参考时标
15.70 h 与 19.02 h（常系数与瞬时首模近似）}
\label{fig:穿越}
\end{figure}

\subsection{敏感性分析}
每组都实测 $\Delta t_f$，见表~\ref{tab:敏感性} 与图~\ref{fig:敏感}。
边界汇取值（$C_e$）是本题\textbf{唯一能改变结论量级的因素}：改用解吸等温线
$C_e=0.1366$ 后 $t_f$ 增大 $22.6$ h（$39.6\%$）。主结果一律采用本文设定的取法，
超出题设的修正只进入敏感性分析并明确标注，其偏移不计入误差分量汇总（7.2 节）。
半径插值两档实测仅差 $0.0022$ h（附件 2 平滑，两插值几乎重合），说明这一取法
对 $t_f$ 的影响很小。

\begin{table}[htbp]
\centering
\caption{敏感性汇总表（问题 3 取法、N=80 对照取法；问题 4 拆分沿用主取法）}
\label{tab:敏感性}
\small
\begin{tabular}{lll}
\toprule
组别 & $t_f$ (h) & $\Delta t_f$ (h) \\
\midrule
$C_e=0.04986$（本文设定，主） & 57.2009 & — \\
$C_e=0.1366$（山药解吸支，超出题设的修正） & 79.8270 & $+22.6261$（$+39.6\%$） \\
环境外推：末值保持（主） & 57.2009 & — \\
环境外推：线性外推（末 1 h 斜率） & 58.8203 & $+1.6194$ \\
环境外推：末 1 h 均值 & 57.5034 & $+0.3025$ \\
半径插值：PCHIP（主，问题 4，附件 2 实测几何对照） & 50.8295 & — \\
半径插值：线性（问题 4） & 50.8317 & $+0.0022$ \\
问题 4 拆分：附录 4 固定 $R$ & 129.1047 & $+78.57$（对内生主解） \\
问题 4 拆分：附录 4 + 内生 $R(t)$（主解） & 50.5369 & — \\
问题 4 拆分：问题 3 & 57.1799 & $+6.64$（对内生主解） \\
\bottomrule
\end{tabular}
\end{table}

\begin{figure}[htbp]
\centering
\includefig[\linewidth]{fig5_敏感性.pdf}{敏感性}
\caption{敏感性条形图：各因素的 $\Delta t_f$ 画成横向条形，其中边界汇取值
$C_e=0.1366$（超出题设的修正）的偏移 $+22.63$ h 最大，环境外推 $+1.62$ h 次之，
半径插值仅 $+0.002$ h}
\label{fig:敏感}
\end{figure}

\subsection{误差与不确定度}

\begin{table}[htbp]
\centering
\caption{误差分量汇总表：按数值误差、环境情景与取法敏感性三类分别报告}
\label{tab:预算}
\small
\begin{tabular}{lp{0.50\linewidth}ll}
\toprule
类别 & 误差源 & 问题 3 (h) & 问题 4 (h) \\
\midrule
数值误差 & 空间网格（Q3：N=160 Richardson 偏差；Q4：连续事件定位 N=80→160 变化） & 0.0164 & 0.0007 \\
数值误差 & 时间步（时间容差收紧十倍） & 0.0014 & 0.0014 \\
数值误差 & 非线性迭代（Picard/Newton 对照） & $1\times10^{-6}$ & $1\times10^{-6}$ \\
数值误差 & 输出重建 / 1 s 插值（曲率上界实测） & $2.6\times10^{-4}$ & $2.6\times10^{-4}$ \\
环境情景 & 环境外推（三档最大偏移） & 1.6194 & 1.6194 \\
对取法的敏感性 & 半径插值（线性与 PCHIP，仅问题 4） & — & 0.0022 \\
对取法的敏感性 & 舍入判定（60 s 网格 + 四位小数，另一种数据使用方式） & 0.0701 & — \\
\midrule
\multicolumn{2}{l}{\textbf{已检验因素累计偏移量级（三角和）}} & $\approx1.71$ & $\approx1.62$ \\
\bottomrule
\end{tabular}
\end{table}

\textbf{模型结构差异}：潜热情景的 $+5.3889$ h 与边界汇取值 $C_e=0.1366$ 的
$+22.6261$ h 单列，不计入上表的合成。前者是 5.3 节在情景假设集合下
无潜热基准的 $t_f$ 偏差量级，后者是超出题设的修正，只作敏感性分析。

\textbf{尚未纳入误差分量汇总的因素}：边界条件的全部可能形式（吸附与水活度
关系缺少数据）、内生收缩标定参数（$k_p$ 的 95\% CI 相对半宽 24\%）、
内部非仿射变形（5.5 节的模型结构分析）。
故问题 3 的达标时间报告为 $t_f=57.18$ h，在所列模型假设与情景集合内，
已检验因素的累计偏移量级约 1.71 h。

\section{模型评价与改进}

\subsection{优点}
\begin{enumerate}
  \item \textbf{统一性与自洽}：四问同为这一套热湿耦合方程、仅取法不同（表~\ref{tab:设定}），
        问题 1 是问题 2 的解耦特例、问题 4 是移动域推广，四问自然衔接；
        质量坐标下收缩几何方程同形、无伪对流项（附录 A 双守恒推导）。
  \item \textbf{这些检验相互独立}：解析谱解与守恒有限体积是两种不同的数学结构，
        V1 给出逐点最大偏差 $5.9\times10^{-7}\,^{\circ}$C 与二阶收敛；此外还有
        退化测试、守恒恒等式（$10^{-11}$ 量级）与参考时标判据检验，分别给出
        可以得到的结论与尚未检验的误差源。
  \item \textbf{收缩由模型自身预测}：问题 4 不读附件 2 而由半经验关系式逐步预测
        $R(t)$（$t_f=50.5369$ h），附件 2 降为验证数据后仍全程吻合
        （0.088 mm）；参数标定与验证用的是同一批数据，标定指标由实测几何
        单向驱动得到，与耦合主解的验证指标是两个不同的量。
  \item \textbf{不确定度的来源逐项列出}：无量纲判据的回代检验给出每条简化的
        $t_f$ 偏移，敏感性分析给出全部 $\Delta t_f$，误差按三类分别报告，
        潜热与边界条件两个情景各自单列。
\end{enumerate}

\subsection{缺点}
\begin{enumerate}
  \item \textbf{湿边界取本文设定的等效边界条件}：$C_{\mathrm{env}}$ 直接作 Robin
        驱动势，缺少吸附等温线的支撑，是本题影响最大的一处假设
        （超出题设的修正 $+22.6261$ h）。
  \item \textbf{无潜热基准}：潜热与显热同量级，基准的 $t_f$ 偏差量级约
        $5.4$ h（表面相变汇情景），更真实的边界条件需要题目未给的吸附与解吸数据。
  \item \textbf{收缩关系式是半经验式}：$g_c,g_p$ 为经验形状函数，力学驱动扩展虽
        更严谨，但在所选幂律保水—简化储能结构内存在精度上限
        （末端 $J=0.3343$，附件 2 的对应值为 0.36），缺少保水曲线与力学实测数据。
  \item \textbf{几何与外推的理想化}：模型取一维径向，端面效应未给出定量界；
        环境条件在 4 h 之后保持末值；内部取等比收缩，内部收缩不可识别。
\end{enumerate}

\subsection{改进}
\begin{enumerate}
  \item 针对缺点 1/2：补测材料吸附/解吸等温线（或水活度关系）与相变焓数据，
        把湿边界与能量边界提升为介质配套的边界条件；
  \item 针对缺点 3：补测真实保水曲线 $p_c(S,T)$ 与骨架模量/松弛谱，
        沿力学驱动路线（5.5 节框架已经建立）重新标定，半经验关系式作为低维验证基准；
  \item 针对缺点 4：补有限长圆柱对照量化端面效应、以更长时程环境数据替换外推、
        在力学框架内放松仿射假设（非仿射已在扩展模型中实现）；
  \item 固定全部比较条件后构造可证明的包络或积分不等式，
        把参考时标提升为已证明的严格下界。
\end{enumerate}

\section*{参考文献}
\begingroup
\renewcommand{\section}[2]{}%
\begin{thebibliography}{9}
\bibitem{luikov} Luikov A V. \emph{Heat and Mass Transfer in Capillary-Porous Bodies}.
Pergamon Press, Oxford, 1966.
\bibitem{crank} Crank J. \emph{The Mathematics of Diffusion}. 2nd ed.,
Oxford University Press, 1975.
\bibitem{pharmacopoeia} 国家药典委员会. \emph{中华人民共和国药典（一部）}.
北京: 中国医药科技出版社, 2020.（药材水分限度的通行标准约 13\%）
\bibitem{saravacos} Saravacos G D. Mass transfer properties of foods. In:
\emph{Engineering Properties of Foods} (M. A. Rao, S. S. Rizvi, eds.),
Marcel Dekker, New York, 1986.
\end{thebibliography}
\endgroup

\appendix
\ctexset{section = {name = {附录 ,：}, number = \Alph{section},
           format = \centering\songti\bfseries\zihao{4}},
         subsection = {number = \Alph{section}.\arabic{subsection},
                       format = \songti\bfseries\zihao{-4}}}

\section{移动域方程的推导要点}
本文的出发点同时包含干物质守恒与水分守恒。设 $\rho_d$ 为单位总体积中的
干物质质量、$v$ 为骨架径向速度、$j$ 为相对骨架的水通量（$j=-\rho_d D\,\partial_r U$）：
\begin{equation}
\partial_t\rho_d+\frac1r\partial_r(r\rho_d v)=0,\qquad
\partial_t(\rho_d U)+\frac1r\partial_r\!\left[r(\rho_d Uv+j)\right]=0 .
\end{equation}
第二式展开后用第一式相减，得物质点形式
\begin{equation}
\rho_d\,(U_t+vU_r)=\frac1r\partial_r\!\left(r\rho_d D\,U_r\right).
\label{eq:物质点}
\end{equation}
引入材料坐标 $q$（干物质标签），物理坐标 $r=\psi(q,t)$，骨架速度
$v=\psi_t|_q$。在\textbf{干物质密度空间均匀}与\textbf{仿射收缩}假设下，取
$\psi=R(t)\sqrt q$，即长度不变、内部等比。附件 2 只识别外半径，
内部收缩不可识别，等比是最简取法。此时 $\rho_d\psi\psi_q$ 与 $t$ 无关，
式~\eqref{eq:物质点} 化为材料坐标形式
$U_t|_q=\partial_q\!\left(4qD/R(t)^2\,U_q\right)$——与固定域同形、无伪对流项。
在移动表面 $q=1$ 处，对表面薄层做水分衡算，几何运动项与通量跳条件中的对应项
恒等相消，Robin 边界 $-D\partial_rU=h_m(U_s-C_{\mathrm{env}})$ 形式不变。温度方程同理。
两点说明：(i) 若由 $U_t|_r=(1/r)\partial_r(rDU_r)$ 出发、假定迁移项相消，所得写法不成立：
该式右侧没有可相消项；上面的材料坐标形式由双守恒方程推出。
(ii) 固定半径退化测试（V5a）\textbf{无法检验}上述推导：
此时骨架速度 $v\equiv0$，移动域与固定域方程逐字相同，
因此该测试只说明两者在数值上一致，不能证明移动域方程本身正确。

\section{参数取值表与收缩参数表}

\begin{table}[H]
\centering
\caption{参数取值表}
\label{tab:参数}
\small
\setlength{\tabcolsep}{5pt}
\begin{tabular}{llll}
\toprule
参数 & 取值/公式 & 取值方式 & 适用 \\
\midrule
$D(C)$ & $7\times10^{-9}\mathrm e^{-0.89/C}$ & 局部 $C$ & 问题 1 \\
$D(C,T)$ & $2.4\times10^{-3}\mathrm e^{-0.45/C}\mathrm e^{-3850/T}$ & 局部 $C$、局部 $T$（开尔文） & 问题 2/3 \\
$D(C,T)$ & $4.2\times10^{-4}\mathrm e^{-0.30/C}\mathrm e^{-3850/T}$ & 局部 $C$、局部 $T$（开尔文） & 问题 4 \\
$h$ & $25$ W/(m$^2\cdot$K) & 常数 & 问题 1--4 \\
$h_m$ & $8\times10^{-7}$ m/s & 常数 & 问题 1--4 \\
$T_a(t),C_{\mathrm{env}}(t)$ & 附件 1（0--14400 s，60 s 步） & 分段线性插值 & 问题 1--4 \\
$R(t)$ & 附件 2（0--259200 s，1800 s 步） & PCHIP 保形插值 & 问题 4 \\
几何 & $L_0=0.25$ m，$R_0=0.02$ m & 一维径向（$L_0/R_0=12.5$） & 问题 1--4 \\
初值 & $U_0=2.55$ kg/kg，$\varTheta_0=301.15$ K & 全域均匀 & 问题 1--4 \\
\midrule
\multicolumn{4}{l}{\textbf{本文设定的假设}} \\
湿边界处理 & $C_{\mathrm{env}}$ 直接作 Robin 驱动势 & 等效边界条件（第三章假设 2） & 问题 1--4 \\
环境外推 & $t>14400$ s 末值保持 & 本文设定（敏感性见 7.1 节） & 问题 2--4 \\
$h,h_m$ 延用 & 附录 2 值延用到问题 2--4 & 本文设定的参数延用假设 & 问题 2--4 \\
半径延拓 & $t>259200$ s 保持 1.198 cm & 本文设定 & 问题 4 \\
\bottomrule
\end{tabular}
\end{table}

\begin{table}[H]
\centering
\caption{内生收缩模型参数表与物理合理性}
\label{tab:收缩参数}
\small
\begin{tabular}{>{\raggedright\arraybackslash}p{0.22\linewidth}lp{0.28\linewidth}p{0.24\linewidth}}
\toprule
参数 & 数值 & 来源 & 物理合理性核对 \\
\midrule
$\varepsilon_c$（塌陷峰值应变） & 0.1108 & 拟合附件 2 & 11.1\% 在毛细压实区间 $[1.4\%,29\%]$ 内 \\
$a,\ b$（塌陷形状指数） & 1.668,\ 1.220 & 拟合附件 2 & 峰值相位 $\bar U\approx1.47$（约 4 h）与急缩段一致 \\
$\varepsilon_p$（停滞幅度） & 0.0709 & 拟合附件 2 & 对应末期孔隙率 30.9\% $\approx$ 附件 2 隐含 31.0\% \\
$k_p$（停滞开启陡峭度） & 3.787 & 拟合附件 2 & 开启相位由 $C_s$ 穿越 $C_{\mathrm{glass}}$ 物理锚定 \\
$C_{\mathrm{glass}}$（玻璃化含水率） & 0.21 & Gordon--Taylor 结构典型参数情景值，未拟合 & 估计区间 0.19--0.23，列入参数不确定清单 \\
$\rho_{sk}$（骨架密度，kg/m$^3$） & 1468 & 生物质骨架密度典型范围值，未拟合 & 结果对其不敏感 \\
\bottomrule
\end{tabular}
\end{table}

\begin{table}[H]
\centering
\caption{量纲检查表}
\label{tab:量纲}
\small
\begin{tabular}{llp{0.68\linewidth}}
\toprule
符号 & 单位 & 核对要点 \\
\midrule
$U,\ C$ & kg/kg（干基） & 边界驱动势 $C_{\mathrm{env}}$ 同基准（本文设定） \\
$D$ & m$^2$/s & $4qD/R^2$ 为 1/s，与 $U_t$ 一致 \\
$T$ & K（不是 ${}^\circ$C） & $\mathrm e^{-3850/T}$ 的 3850 单位为 K；读入与全程统一取开尔文并作单位校验 \\
$h,h_m$ & W/(m$^2\cdot$K)、m/s & $-k\partial_rT=h(T_s-T_a)$ 两侧 W/m$^2$；$h_m\Delta U$ 为 m/s \\
$a_{\mathrm{eff}}$ & J/(m$^3\cdot$K) & $\rho c_p$ 逐点取值，$\rho$ 不进湿分方程 \\
$q$ & — & $q=(X/R_0)^2\in[0,1]$，干物质标签 \\
$\varepsilon_c,\varepsilon_p$ & —（应变） & 式~\eqref{eq:收缩闭合} 的乘子修正 \\
\bottomrule
\end{tabular}
\end{table}

\section{内生收缩模型的标定细节}
内生收缩模型（式~\eqref{eq:收缩闭合}）的 5 个参数
$\varepsilon_c,a,b,\varepsilon_p,k_p$ 由 scipy 的 \texttt{least\_squares}
在附件 2 的 145 个采样点上对 $R_{\mathrm{pred}}-R_{\mathrm{data}}$ 做带物理边界约束的
最小二乘估计。$C_{\mathrm{glass}}=0.21$ 取 Gordon--Taylor 结构典型参数情景值，
$\rho_{sk}=1468$ kg/m$^3$ 取生物质骨架密度典型范围值，两者不参与拟合，
均列入参数不确定清单。
过拟合检验用奇偶折半留出法：两折上训练与留出 RMSE 均在 0.045--0.050 mm，
且两折参数互差 $<10\%$。该检验只检验了同一光滑曲线上的插值，
因为湿度场由附件 2 实测几何生成，参数标定与验证用的是同一批数据，
所以只能说明插值稳健，不构成参数被独立识别的证明；
块 bootstrap（残差自相关块长 24、$B=200$）95\%
置信区间的相对半宽 $\varepsilon_c$ 4.7\%、$a$ 11.9\%、$b$ 11.0\%、
$\varepsilon_p$ 3.8\%、$k_p$ 23.9\%，分布有偏斜或多峰迹象；
变初值 $\times$ 变损失函数的最大参数偏差 $\le0.98\%$；
各机制在其主导区段单独拟合与全曲线拟合的幅度参数差 $<2\%$，独立识别需完整耦合标定。

\textbf{参数的第二路线检验}：$\varepsilon_c=11.1\%$ 落在毛细压
$2\gamma/r_{\mathrm{pore}}$ 对基体模量的理论应变区间 $[1.4\%,28.8\%]$；
$\varepsilon_p$ 对应的末期孔隙率 30.9\% 与由“平台半径和质量守恒”独立推出的
31.0\% 一致；$C_{\mathrm{glass}}=0.21$ 与“附件 2 冻结起点约 21 h
$\leftrightarrow$ 模型 $\bar U(21\ \mathrm{h})=0.234$”交叉吻合；
$\rho_{sk}=1468$ 在生物质骨架密度典型范围内且结果对其不敏感。
第二路线检验的结果是：$\varepsilon_c$、$a$、$b$、$\varepsilon_p$、$C_{\mathrm{glass}}$、
$\rho_{sk}$ 均获得支持；\textbf{$k_p$ 存疑}：它只有曲线拟合一条路，
没有独立的物理观测量，置信区间也最宽，达 24\%；但 $k_p$ 只影响过渡带形状，
不触及达标时间与平台半径的结论。
力学扩展模型的标定（4 起点同一解、留出、bootstrap、末端平衡分析）
见 5.5 节正文与数据文件 \texttt{03-数据/mechRefit.csv}。

\section{复现说明与代码结构}
一条命令 \texttt{python run\_all.py} 依次执行下列脚本；Windows 下需先设置
环境变量 \texttt{PYTHONIOENCODING=utf-8}。全部脚本只用到
numpy/scipy/pandas/openpyxl。
\begin{itemize}
  \item \texttt{run\_q1.py}：问题 1 求解 + V1 解析比较 + 守恒与退化检验 + result1.xlsx；
  \item \texttt{run\_q23.py}：问题 2/3 耦合求解 + $t_f$ + V3/V4/V6 判据检验 +
        result2/result3.xlsx + 环境外推三档；
  \item \texttt{run\_q4.py}：问题 4 移动域 + 效应拆分 + V5 退化 + result4.xlsx；
        \texttt{run\_q4\_endo.py}：内生几何主解（满 72 h 积分、三段误差评价）；
  \item \texttt{run\_m4.py}：无量纲判据回代检验 + 敏感性汇总 + 误差分量汇总 +
        V2 尾段 + 结果汇总；
  \item \texttt{latent\_scenario.py}：潜热对照情景（5.3 节）；
  \item \texttt{run\_mech\_refit.py}（03-数据）：力学扩展模型重标定（5.5 节）。
\end{itemize}
全部结果文件确定性写出（固定归档时间戳），连续两次运行 result1--4 的 md5
完全一致（result1 \texttt{22eb47f7…}、result2 \texttt{12cde4e8…}、
result3 \texttt{fbe982ff…}、result4 \texttt{80c73431…}）。

\subsection*{支撑材料文件列表}
交付文件夹为 \texttt{论文/}（顶层目录，63 个文件、7 个目录，未压缩约 5.0 MB）；
四个结果工作簿按题目要求命名并置于包根目录，其余材料在 \texttt{论文/支撑材料/} 下
按源码仓库结构分组。下列清单与该交付文件夹内容逐项一致。

\textbf{如何复现}：打开（或解包）后进入 \texttt{论文/支撑材料/}，该目录完整镜像源码仓库的相对结构
（\texttt{02-代码/}、\texttt{03-数据/}、\texttt{04-结果/}、\texttt{05-图表/}、\texttt{06-论文/}），
脚本按仓库相对路径解析的输入输出路径、以及本文的插图与源码引用路径均\textbf{原样成立}，无需调整目录。
编译论文：\texttt{cd 支撑材料/06-论文 \&\& xelatex 中药材烘干热湿耦合建模、守恒求解与模型检验.tex}
（连跑两遍以解析交叉引用）；复现结果：\texttt{cd 支撑材料/02-代码 \&\& python run\_all.py}，
按既定顺序执行全部脚本并由脚本写出四个结果工作簿与 \texttt{03-数据/} 的中间数据
（\texttt{03-数据/} 下的 \texttt{.npz} 张量缓存合计约 23.5 MB，超 20 MB 上限，由该命令重新生成，故未随包提供）。
\texttt{04-结果/} 是脚本的结果写入处，交付结果工作簿同时置于包根目录，故该目录在包内留空，
保留它是为了复跑时不必手工建目录。赛题原件按规范第十一条豁免未随包。

{\raggedright
\newcommand{\fsep}{、\allowbreak}
\begin{itemize}
  \item \textbf{论文/}（6）：\texttt{中药材烘干热湿耦合建模、守恒求解与模型检验.pdf}\fsep
        \texttt{result1.xlsx}\fsep \texttt{result2.xlsx}\fsep \texttt{result3.xlsx}\fsep
        \texttt{result4.xlsx}\fsep \texttt{AI工具使用声明.docx}；
  \item \textbf{论文/支撑材料/02-代码/}（13）：\texttt{analytic\_heat.py}\fsep
        \texttt{endogenous\_calibrated.py}\fsep \texttt{latent\_scenario.py}\fsep
        \texttt{mech\_shrinkage.py}\fsep \texttt{plot\_all.py}\fsep \texttt{run\_all.py}\fsep
        \texttt{run\_m4.py}\fsep \texttt{run\_q1.py}\fsep \texttt{run\_q23.py}\fsep
        \texttt{run\_q4.py}\fsep \texttt{solver\_q1.py}\fsep \texttt{solver\_q23.py}\fsep
        \texttt{solver\_q4.py}；
  \item \textbf{论文/支撑材料/03-数据/}（20）：\texttt{run\_mech\_A.py}\fsep
        \texttt{run\_mech\_mainline.py}\fsep \texttt{run\_mech\_refit.py}\fsep
        \texttt{run\_q4\_endo.py}\fsep \texttt{run\_q4\_endo\_fields.py}\fsep
        \texttt{交付口径-核对.py}\fsep \texttt{结果总账.csv}\fsep \texttt{paper\_tables.csv}\fsep
        \texttt{q4\_endo\_eval.csv}\fsep \texttt{q4\_split.csv}\fsep
        \texttt{q4\_endo\_convergence.csv}\fsep \texttt{mechRefit.csv}\fsep
        \texttt{sensitivity.csv}\fsep \texttt{error\_budget.csv}\fsep
        \texttt{latent\_scenario.csv}\fsep \texttt{D\_of\_C.csv}\fsep \texttt{surface\_flux.csv}\fsep
        \texttt{交付口径.md}\fsep \texttt{力学升级\_修复重标定.md}\fsep \texttt{内生收缩模型.md}；
  \item \textbf{论文/支撑材料/04-结果/}（0）：空目录（脚本的结果写入处，交付结果同时在包根目录）；
  \item \textbf{论文/支撑材料/05-图表/}（17）：\texttt{fig1\_输入数据.pdf}\fsep
        \texttt{fig2\_问题1解析验证.pdf}\fsep \texttt{fig3\_温湿剖面.pdf}\fsep
        \texttt{fig4\_达标穿越.pdf}\fsep \texttt{fig5\_敏感性.pdf}\fsep
        \texttt{fig6\_收缩机制分析.pdf}\fsep \texttt{fig7\_内生收缩贴合.pdf}\fsep
        \texttt{fig9\_截面演化.pdf}\fsep \texttt{fig10\_力学升级.pdf}\fsep
        \texttt{fig11\_扩散系数与时间尺度.pdf}\fsep \texttt{fig12\_温度场时空云图.pdf}\fsep
        \texttt{fig13\_水分场时空云图.pdf}\fsep \texttt{fig14\_表面通量与表面扩散系数.pdf}\fsep
        \texttt{fig15\_收敛与守恒检验.pdf}\fsep \texttt{fig16\_中心表面演化与半径收缩.pdf}\fsep
        \texttt{fig17\_几何模型与边界条件示意图.pdf}\fsep
        \texttt{\_fig4\_对照\_附4固定R\_maxU.csv}；
  \item \textbf{论文/支撑材料/06-论文/}（7）：\texttt{中药材烘干热湿耦合建模、守恒求解与模型检验.tex}\fsep
        \texttt{tables.tex}\fsep \texttt{gen\_tables.py}\fsep
        \texttt{check\_abstract.py}\fsep \texttt{check\_numbers.py}\fsep
        \texttt{check\_refs.py}\fsep \texttt{check\_table6.py}。
\end{itemize}}

{\emergencystretch=2em
其中 \texttt{result1.xlsx}、\texttt{result3.xlsx}、\texttt{result4.xlsx} 为原表原样入包；
包根目录的 \texttt{result2.xlsx} 为\textbf{降采样版}（仓库内文件名为 \texttt{result2\_降采样.xlsx}）：
原表为 1 s 时间网格（温度与水分浓度 2 张表 × 209447 数据行 × 22 列，27.9 MB），
超过支撑材料 20 MB 上限，故按每 1800 s（30 min）一帧抽取、每帧保留完整径向剖面（21 个节点），
并加留首帧 1 s 与末帧 209447 s（$t_f+3600$ s 余量），共 118 帧；
表名、列头、列序与原表一致。其余中间工作簿与张量（\texttt{.npz}）、位图、日志仍不入包。}

核心求解步骤的程序摘录（守恒步进与表面重建，注释从略）如下。

\begin{lstlisting}[language=Python,caption={守恒耦合步进核心（solver 摘录）}]
# 表面值：半单元内阻与外对流阻串联（D_eff 取半单元中点值）
def surface_m(CN, Ts, Cenv, dq):
    g = lambda Cs: 8.0*D_of(0.5*(Cs+CN), Ts)*(Cs-CN)/(R0**2*dq) \
                   - 2.0*HM*(Cenv-Cs)/R0
    Cs = brentq(g, lo, hi, xtol=1e-15, rtol=1e-14)
    return Cs, 2.0*HM*(Cenv-Cs)/R0

# 一个耦合后向欧拉步：热场（冻结系数，线性）→ 湿场（冻结温度，Newton）
def coupled_step(U, T, dt, Ta, Cenv, dq, newton):
    ch = 4.0*qf*k_of(Uf)*dt/(R0**2*dq**2)      # 面通量系数
    Rh = R0**2*dq/(8.0*ks) + R0/(2.0*H)        # 表面串联热阻
    T_new = solve_banded((1,1), ab_T, rhs_T)   # 热场线性求解
    for _ in range(30):                        # 湿场 Newton/Picard
        dU = solve_banded((1,1), J, -G)
        U_new = U_new + dU
    return U_new, T_new
\end{lstlisting}

\section{核心程序代码}
本附录列出四个核心模块的完整源码，共 868 行；四份源码均自源文件直接引入，未手工转录。
其余脚本、数据与图表的完整可运行代码见支撑材料 \texttt{论文/支撑材料/02-代码/}（13 个脚本）
与 \texttt{论文/支撑材料/03-数据/}（6 个脚本与全部证据数据）。

\subsection{解析谱解}
圆柱 Robin 特征展开与时变环境 Duhamel 卷积的解析解，用作问题 1 的独立参照（V1）。
\lstinputlisting[language=Python,caption={解析谱解（analytic\_heat.py，122 行）}]{../02-代码/analytic_heat.py}

\subsection{热湿耦合主求解器}
质量坐标下的守恒型有限体积格式：半单元内阻与外对流阻串联重建表面值、
后向欧拉加局部外推、Picard--Newton 交替迭代、自适应步长。
\lstinputlisting[language=Python,caption={热湿耦合主求解器（solver\_q23.py，271 行）}]{../02-代码/solver_q23.py}

\subsection{移动域求解器}
仿射收缩 $r=R(t)\sqrt q$ 下的移动域求解与达标判定。
\lstinputlisting[language=Python,caption={移动域求解器（solver\_q4.py，232 行）}]{../02-代码/solver_q4.py}

\subsection{内生收缩闭合}
由湿度场逐步预测半径 $R(t)$ 的半经验闭合（三项机制），标定参数见附录 C。
\lstinputlisting[language=Python,caption={内生收缩闭合（endogenous\_calibrated.py，243 行）}]{../02-代码/endogenous_calibrated.py}

\clearpage
\section*{AI 工具使用声明}
\markboth{AI 工具使用声明}{AI 工具使用声明}

本参赛队就本次竞赛中 AI 工具的使用情况声明如下：

\begin{table}[H]
\centering
\small
\begin{tabular}{@{}cp{0.15\linewidth}p{0.33\linewidth}p{0.22\linewidth}c@{}}
\toprule
序号 & AI 工具名称 & 使用场景/用途 & 涉及论文章节 & 使用程度 \\
\midrule
1 & KimiCode、DeepSeek & 程序测试：辅助编写与调试数值求解代码，以及守恒恒等式、网格与时间收敛、退化测试等检验脚本 & 第五章（模型建立与求解）、第六章（模型检验）、附录 D & 中度 \\
2 & KimiCode、DeepSeek & 绘图：辅助整理数据并生成全文插图 & 第一章至第七章的插图 & 中度 \\
3 & KimiCode、DeepSeek & 数学公式推理验证：辅助核对控制方程与移动域推导、解析谱解、量纲与极限退化 & 第 5.1 节、附录 A、第 6.1--6.2 节 & 轻度 \\
\bottomrule
\end{tabular}
\end{table}

\textbf{使用程度说明：}

轻度：仅用于语法检查、文字润色、代码片段解释、公式验证等辅助性工作，不影响论文核心内容与结论。

中度：用于辅助代码编写、文献整理、思路拓展等，核心建模与分析仍由队员独立完成。

本队承诺以上声明内容全部属实，论文核心建模方法、算法设计、结果分析与结论均为队员独立完成，不存在 AI 工具直接生成核心内容或代写代做的情况。如有不实，愿承担全部责任并接受相应处理。

\vspace{1em}
\noindent 全体队员签字：\underline{\hspace{4cm}} \hfill 日期：\underline{\hspace{3cm}}

\end{document}
